% Ch11 WS1 -- solution composition. Block 1. CED content.
\documentclass{apchem-worksheet}

\apcunitword{Chapter}
\apcunit{11}
\apcunittitle{Properties of Solutions}
\apcdoctitle{Worksheet 1 \textbullet\ Solution Composition}
\apcsubtitle{Zumdahl \S11.1 \textbullet\ litres of SOLUTION vs kilograms of SOLVENT}

\begin{document}
\wsheader[State which denominator each calculation uses before you compute.
Most errors on this sheet are denominator errors, not arithmetic errors.]

\problem[]{Write the defining formula for each:}
\begin{parts}
  \item Molarity: \answerblank[6cm]{mol solute / L solution}
  \item Molality: \answerblank[6cm]{mol solute / kg solvent}
  \item Mass percent: \answerblank[7.0cm]{(mass solute / mass solution)$\times$100}
  \item Mole fraction: \answerblank[6cm]{$n_A / n_{\text{total}}$}
\end{parts}

\problem[]{A solution is made by dissolving \SI{20.0}{\gram} of \ce{KCl}
($M = 74.55$) in \SI{180.0}{\gram} of water. The final volume is
\SI{190.}{\milli\liter}.}
\begin{parts}
  \item Moles of \ce{KCl}: \answerblank[3cm]{\SI{0.268}{\mole}}
  \item Molarity: \answerblank[3cm]{\SI{1.41}{\Molar}}
  \item Molality: \answerblank[3cm]{\SI{1.49}{\mole\per\kilo\gram}}
  \item Mass percent: \answerblank[3cm]{10.0\%}
  \item Mole fraction of \ce{KCl}: \workspace[2cm]
        \keyonly{\begin{apcanswer}$n(\ce{H2O}) = 180.0/18.02 =
        \SI{9.989}{\mole}$; $\chi = 0.268/(0.268+9.989) =
        0.0261$\end{apcanswer}}
\end{parts}

\problem[]{Explain why molarity changes with temperature but molality does
not. \answerlines[3]{Molarity is per litre of solution, and volume expands
when heated, so the same moles occupy more litres and the molarity falls.
Molality is per kilogram of solvent, and mass is unaffected by temperature,
so it stays constant.}}

\problem[]{A student adds \SI{0.500}{\mole} of solute to exactly
\SI{1.000}{\liter} of water and labels the bottle
``\SI{0.500}{\Molar}.''}
\begin{parts}
  \item Identify the error. \answerlines[2]{Molarity requires litres of
        \emph{solution}. Adding solute increases the total volume beyond
        1.000 L, so the concentration is less than 0.500 M.}
  \item What quantity \emph{is} exactly 0.500 for this solution?
        \answerblank[5.6cm]{the molality (per kg water)}
  \item How should the solution have been prepared?
        \answerlines[2]{Dissolve the solute in less than a litre of water,
        then add water up to the 1.000 L calibration mark of a volumetric
        flask.}
\end{parts}

\problem[]{Which concentration measure is most appropriate, and why?}
\begin{parts}
  \item Calculating the partial pressure of one gas in a mixture:
        \answerblank[5cm]{mole fraction}
  \item Titrating an acid of known volume:
        \answerblank[5cm]{molarity}
  \item Reporting the composition of a saline solution to a patient:
        \answerblank[5cm]{mass percent}
\end{parts}

\problem[]{A \SI{500.0}{\gram} sample of aqueous glucose
(\ce{C6H12O6}, $M = 180.16$) is 12.0\% glucose by mass.}
\begin{parts}
  \item Mass of glucose present: \answerblank[3cm]{\SI{60.0}{\gram}}
  \item Moles of glucose: \answerblank[3cm]{\SI{0.333}{\mole}}
  \item Mass of water: \answerblank[3cm]{\SI{440.0}{\gram}}
  \item Molality: \answerblank[3cm]{\SI{0.757}{\mole\per\kilo\gram}}
\end{parts}

\begin{spiralstrip}{Chapters 9--10 \textbullet\ spiral}
  \item Hybridization of C in \ce{CH4}: \answerblank[2cm]{sp$^3$}
  \item Why does ice float? \answerblank[4.6cm]{open H-bonded lattice}
  \item $\Delta H_{\text{vap}}$ of water:
        \answerblank[3.2cm]{\SI{40.7}{\kilo\joule\per\mole}}
  \item Type of solid: graphite? \answerblank[3.4cm]{network covalent}
  \item Which bond type blocks rotation? \answerblank[1.8cm]{$\pi$}
\end{spiralstrip}

\end{document}
